\chapter*{Bibliographie}
\addcontentsline{toc}{chapter}{Bibliographie}
\begin{itemize}\small
\item Bargain O., Vicard A. (2014), « Le RMI et son successeur le RSA découragent-ils certains jeunes de travailler ? Une analyse sur les jeunes autour de 25 ans », \emph{Économie et Statistique}.
\item Beer S., de Mooij R., Liu L. (2020), « International corporate tax avoidance: a review of the channels, magnitudes, and blind spots », \emph{Journal of Economic Surveys}.
\item Benzarti Y., Carloni D., Harju J., Kosonen T. (2020), « What goes up may not come down: asymmetric incidence of value-added taxes », \emph{Journal of Political Economy}.
\item Blundell R., Pistaferri L., Preston I. (2008), « Consumption inequality and partial insurance », \emph{American Economic Review}.
\item Boissel C., Matray A. (2022), « Dividend taxes and the allocation of capital », \emph{American Economic Review}.
\item Bozio A., Breda T., Grenet J. (2019), « Does tax-benefit linkage matter for the incidence of social security contributions? », document de travail IZA.
\item Cabannes P.-Y., Houdré C., Landais C. (2014), « Comment le revenu imposable des ménages répond-il aux variations des taux marginaux supérieurs ? », \emph{Revue économique}.
\item Carbonnier C. (2007), « Who pays sales taxes? Evidence from French VAT reforms, 1987-1999 », \emph{Journal of Public Economics}.
\item Carbonnier C., Malgouyres C., Py L., Urvoy C. (2022), « Who benefits from tax incentives? The heterogeneous wage incidence of a tax credit », \emph{Journal of Public Economics} (document de travail 2019).
\item Chirinko R., Fazzari S., Meyer A. (1999), « How responsive is business capital formation to its user cost? », \emph{Journal of Public Economics}.
\item CNEPI — France Stratégie (2021), \emph{Évaluation du crédit d'impôt recherche}.
\item Commission européenne, base AMECO (écart de production, production potentielle), via DBnomics.
\item Dechezleprêtre A. et al., travaux de l'OCDE sur les effets des prix de l'énergie sur l'emploi manufacturier.
\item Dussaux D. (2020), « The joint effects of energy prices and carbon taxes on environmental and economic performance: evidence from the French manufacturing sector », OCDE.
\item FMI (2014), « Is it time for an infrastructure push? The macroeconomic effects of public investment », \emph{World Economic Outlook}, chapitre 3.
\item FMI, \emph{World Economic Outlook}, écart de production (NGAP\_NPGDP), via DBnomics.
\item France Stratégie (2020), \emph{Évaluation du crédit d'impôt pour la compétitivité et l'emploi}, rapport du comité de suivi.
\item Fuest C., Peichl A., Siegloch S. (2018), « Do higher corporate taxes reduce wages? Micro evidence from Germany », \emph{American Economic Review}.
\item Hall B., Mairesse J., Mohnen P. (2010), « Measuring the returns to R\&D », \emph{Handbook of the Economics of Innovation}.
\item Havranek T. (2015), « Measuring intertemporal substitution: the importance of method choices and selective reporting », \emph{Journal of the European Economic Association}.
\item Heckemeyer J., Overesch M. (2017), « Multinationals' profit response to tax differentials », \emph{Canadian Journal of Economics}.
\item Le Barbanchon T., Rathelot R., Roulet A. (2017), « Unemployment insurance and reservation wages: evidence from administrative data », \emph{Journal of Public Economics}.
\item Marin G., Vona F. (2021), « The impact of energy prices on socioeconomic and environmental performance: evidence from French manufacturing establishments », \emph{European Economic Review}.
\item Merlo V., Wamser G., estimations récentes du transfert de bénéfices des multinationales.
\item OFCE, DARES : revues des élasticités d'offre de travail (participation et heures).
\item Ramey V. (2019), « Ten years after the financial crisis: what have we learned from the renaissance in fiscal research? », \emph{Journal of Economic Perspectives}.
\item Rapport v0.6 du projet (18 novembre 2026), \emph{L'économie française au banc d'essai}.
\item Rapport v0.7 du projet (23 novembre 2026), \emph{Comportements économiques, fiabilisation du modèle et première optimisation}.
\item Schmieder J., von Wachter T. (2016), « The effects of unemployment insurance benefits: new evidence and interpretation », \emph{Annual Review of Economics}.
\item Zwick E., Mahon J. (2017), « Tax policy and heterogeneous investment behavior », \emph{American Economic Review}.
\end{itemize}
